\chapter[DPF for HMC in DSGE and MacroFinance]{Assessing Differentiable Particle Filters for HMC in Nonlinear DSGE and MacroFinance Models}
\label{ch:bf-dpf-hmc-assessment}

This chapter asks the program-level question left open by the earlier survey and
implementation chapters: is the current differentiable particle-filter path a
justified route to Hamiltonian Monte Carlo for nonlinear DSGE and MacroFinance
models?

The answer from the current literature is neither a simple yes nor a simple no.
The literature and the student baselines support differentiable particle filters
as a serious research direction for nonlinear filtering and for
gradient-informed parameter inference.  They do \emph{not} yet eliminate the
need to check, rung by rung, what scalar is being differentiated, whether that
scalar is an exact likelihood, an unbiased estimator, or a biased surrogate,
and whether the resulting value-gradient path is stable enough for HMC in the
structural model classes that matter to BayesFilter.

\section{The HMC question is a target question}
\label{sec:bf-dpf-hmc-question}

BayesFilter's HMC doctrine in Chapter~\ref{ch:bf-hmc-for-state-space} is that an
HMC transition is valid only relative to a declared target.  A filter that
returns a value is not automatically an HMC backend; a filter that returns a
differentiable value is still not automatically an HMC backend.  The decisive
question is whether the gradient used by the integrator corresponds to the same
log target used in the Metropolis correction.

In the DPF context, this target question becomes particularly sharp.  The
algorithm may involve:
\begin{itemize}
  \item a deterministic or stochastic flow approximation;
  \item an importance-weight correction;
  \item a relaxed resampling map;
  \item an amortized neural surrogate for the relaxed resampling map;
  \item frozen or evolving random numbers across leapfrog steps.
\end{itemize}
Each of these layers can change the object being evaluated or differentiated.
The current literature contains valid constructions for some of these layers,
but it does not collapse them into a universal theorem that any differentiable
particle filter yields a valid HMC target.

\section{What the literature supports by rung}
\label{sec:bf-dpf-rung-assessment}

The most useful way to assess suitability is rung by rung.

\subsection{EDH alone}
\label{sec:bf-dpf-hmc-edh}

EDH is strongest as a filtering and derivation benchmark.  In the
linear-Gaussian special case, it provides a valuable exact-recovery test against
Kalman filtering, and in broader nonlinear settings it provides a deterministic
transport approximation motivated by a continuous homotopy.  The literature and
CIP derivations support EDH as a meaningful proposal or approximation tool.

That is not the same as saying EDH alone defines the correct HMC target for a
nonlinear DSGE or MacroFinance model.  Once the global Gaussian closure ceases
to be exact, EDH becomes a controlled approximation.  It may still be a useful
value-side likelihood surrogate or a precursor to PF-PF correction, but the
literature does not justify treating raw EDH as a drop-in exact HMC likelihood
path for the target models BayesFilter cares about.

\subsection{EDH/PF with importance weights}
\label{sec:bf-dpf-hmc-edhpf}

This is the first rung at which a clearer target story emerges.  By using the
flow as a proposal and then correcting with the change-of-variables weight,
PF-PF retains a standard importance-sampling interpretation.  For BayesFilter,
this makes EDH/PF the first rung that is plausibly aligned with HMC-compatible
value-side development.

Plausibly aligned is still not the same as ready.  Three further questions
remain:
\begin{enumerate}[label=(\roman*)]
  \item Is the estimator variance small enough that repeated HMC evaluations are
    meaningful in practice?
  \item Is the Jacobian or log-determinant path stable enough for automatic
    differentiation and compiled execution?
  \item Is the resulting value-gradient path finite and reproducible in the
    structural models of interest?
\end{enumerate}
The literature motivates these questions but does not close them for nonlinear
DSGE or MacroFinance systems.

\subsection{Soft or other smooth resampling relaxations}
\label{sec:bf-dpf-hmc-soft}

Soft resampling is attractive because it yields a smoother pathwise derivative.
The literature and the CIP resampling chapter also make the main caveat plain:
it is a relaxation that changes the resampling map.  This means it is naturally
interpreted as a biased differentiable surrogate unless a more specific theorem
is supplied.

For BayesFilter's HMC program, that caveat matters a great deal.  HMC can still
be useful on a biased differentiable surrogate if the project consciously
chooses that surrogate as the target-defining object or uses it for diagnostic
or optimization stages.  What BayesFilter should not do is quietly promote the
surrogate posterior to an exact posterior claim.

The student \texttt{advanced\_particle\_filter} notebooks are helpful here.  By
running linear-Gaussian-reference diagnostics and then HMC experiments on a
stochastic-volatility state-space model, they effectively show that resampling
choice can visibly affect posterior summaries.  That is exactly the evidence one
would expect if the resampling mechanism is part of the target definition rather
than an invisible computational convenience.

\subsection{OT resampling}
\label{sec:bf-dpf-hmc-ot}

Entropy-regularized OT resampling is more structured than simple soft
resampling.  It has a clearer geometric interpretation and a more direct
literature base through \citet{corenflos2021differentiable}.  For BayesFilter,
this makes OT resampling the strongest currently available differentiable
resampling candidate.

Even here, however, the literature supports a careful formulation rather than a
free upgrade.  The transport is regularized, finite-iteration Sinkhorn solves
introduce their own approximation layer, and the barycentric projection is not
identical to categorical multinomial resampling.  The advanced student repo's
sensitivity sweeps are therefore directly relevant to BayesFilter's HMC
decision.  If posterior summaries are sensitive to the transport regularization
or solver budget, then the resampling layer is not just a numerical backend; it
is part of the target specification.

\subsection{LEDH plus neural OT}
\label{sec:bf-dpf-hmc-ledh-neuralot}

This is the intended final target of the program, not the currently justified
starting point.  The literature supports the constituent ideas: local particle
flow, transport-based differentiable resampling, and learned acceleration of OT
maps.  The student \texttt{advanced\_particle\_filter} repository is especially
useful here because it exposes the practical promise of the combination: large
runtime speedups and a clear experiment path from Sinkhorn to amortized OT.

For nonlinear DSGE and MacroFinance HMC, however, this rung also carries the
largest stack of assumptions:
\begin{enumerate}[label=(\roman*)]
  \item local linearization adequacy in the flow;
  \item stability of the local Jacobian and log-determinant path;
  \item approximation error from entropic OT regularization;
  \item approximation error from the learned OT operator;
  \item training-distribution mismatch when the learned operator is used on new
    model classes or state dimensions.
\end{enumerate}
This makes LEDH + neural OT a promising research destination, but not yet the
best justified first production DPF target for BayesFilter.

\section{Why nonlinear DSGE models are especially hard}
\label{sec:bf-dpf-dsge-hard}

Nonlinear DSGE models impose structural constraints that are not well described
by small generic filtering benchmarks alone.  Chapter~\ref{ch:bf-structural-
  deterministic-dynamics} showed that mixed structural models are governed by a
pushforward law over declared stochastic coordinates and deterministic
completion coordinates.  This already creates a subtle contract issue for
Gaussian approximations and sigma-point methods.  It is equally important for
particle methods.

A DPF backend for nonlinear DSGE models must cope with:
\begin{itemize}
  \item structural state partitions rather than homogeneous latent vectors;
  \item parameter regions where the model solution or observation map may become
    invalid or only barely valid;
  \item likelihood surfaces whose geometry is shaped by perturbation
    approximations, determinacy constraints, or pruning conventions;
  \item potentially stiff local observation Jacobians when measurements are
    informative or nearly deterministic.
\end{itemize}
These features do not rule out DPFs.  They do imply that filtering quality on
small generic SSMs is not enough.  The first HMC-relevant question is whether
the filter respects the structural target contract at all; the second is whether
its derivative path remains faithful and finite in the intended parameter
region.

The current literature on DPFs is much richer for geophysical and generic
nonlinear tracking models than for nonlinear DSGE likelihood evaluation.  That
means BayesFilter is not merely implementing an established macroeconomic
standard.  It is carrying frontier filtering ideas into a structural-economics
setting whose target and derivative requirements are unusually strict.

\section{Why MacroFinance models are also hard}
\label{sec:bf-dpf-macrofinance-hard}

MacroFinance latent-factor models present a different but equally demanding
regime.  Their challenges include:
\begin{itemize}
  \item long observation horizons;
  \item missing-data and mixed-frequency structure;
  \item latent-factor dimensions large enough for particle degeneracy to matter
    materially;
  \item sharp likelihood curvature in volatility and measurement-noise
    directions;
  \item HMC chains that are sensitive to compiled-path failures and finite-value
    boundary behavior.
\end{itemize}
Chapter~\ref{ch:bf-nawm-design-target} explains why industrial-scale HMC should
prefer analytic or custom derivative paths where available.  DPFs do not evade
that concern; they intensify it.  A particle method can help represent
non-Gaussian filtering structure, but it may also introduce Monte Carlo
variability, transport approximation bias, and resampling-surrogate bias into
the same log target that HMC is trying to integrate accurately.

This does not make DPFs irrelevant to MacroFinance.  On the contrary, the
non-Gaussian and nonlinear features of large latent-factor models are exactly
why DPFs are interesting.  It does mean that BayesFilter should treat DPFs as a
candidate track to be benchmarked against exact or controlled-approximate
likelihood paths rather than as an automatic replacement for them.

\section{Comparison against the student conclusions}
\label{sec:bf-dpf-hmc-student-comparison}

The student materials are useful precisely because they are more ambitious than
BayesFilter's current doctrine in some places.

The MLCOE report frames particle flow and differentiable particle filtering as a
potential route toward high-dimensional financial applications and repeatedly
emphasizes TensorFlow implementation and HMC compatibility.  This is valuable as
a research motivation, and it surfaces literature and benchmark ideas that would
otherwise be missing from the BayesFilter monograph.  Where BayesFilter should
be more careful is in the inferential step from promising filtering technology
to justified HMC target technology.  The report's application discussion is more
forward-looking than source-backed for nonlinear DSGE or MacroFinance likelihood
contracts.

The \texttt{advanced\_particle\_filter} repository goes further by running HMC
experiments with soft, Sinkhorn, and amortized OT resamplers on a
stochastic-volatility SSM.  This is useful because it shows what the empirical
questions look like when one actually tries to sample with these targets.  It is
also a reason for caution: the repository itself is explicit that the amortized
operator produces a measurable approximation residual and a visible posterior
shift relative to reference comparisons.  BayesFilter should treat this as
evidence that the target question remains active, not as evidence that the
surrogate can be regarded as harmless.

Thus BayesFilter's more conservative conclusion is not that the student work is
wrong.  It is that the student work is experimenting on the frontier, while
BayesFilter's doctrine is trying to define what would have to be true before the
frontier method can be trusted inside a reusable HMC likelihood stack for
structural models.

\section{Assessment table}
\label{sec:bf-dpf-assessment-table}

Table~\ref{tab:bf-dpf-hmc-assessment} summarizes the current assessment.

\begin{table}[ht]
\centering
\small
\begin{tabular}{p{0.18\textwidth} p{0.19\textwidth} p{0.16\textwidth} p{0.28\textwidth}}
\hline
Rung & What the literature supports & Current HMC status & Main missing evidence \\
\hline
EDH & Exact linear-Gaussian recovery and useful nonlinear transport approximation & Benchmark / value-side only & Nonlinear target-validity story beyond approximation \\
EDH/PF with weights & Proposal-correction path with principled weights & First serious HMC candidate & Variance, Jacobian stability, compiled-path evidence \\
Soft differentiable resampling & Smooth surrogate gradients & Surrogate only unless target is declared approximate & Bias quantification and target interpretation \\
OT resampling & Strongest current differentiable-resampling literature base & Plausible surrogate or approximate-target candidate & Sensitivity to regularization and solver budget \\
LEDH + neural OT & Promising frontier research target & Not first production rung & Local-flow bias + learned-operator bias + transfer evidence \\
\hline
\end{tabular}
\caption{Current assessment of the differentiable particle-filter ladder for HMC use in BayesFilter.}
\label{tab:bf-dpf-hmc-assessment}
\end{table}

\section{Program-level recommendation}
\label{sec:bf-dpf-program-recommendation}

The literature survey, the student comparisons, and BayesFilter's current HMC
contract together support the following recommendation.

\begin{enumerate}[label=(\roman*)]
  \item BayesFilter should continue the DPF program.  The path is technically
    serious and the literature supports it as an important frontier for
    nonlinear filtering and gradient-informed parameter inference.
  \item BayesFilter should not skip directly to LEDH plus neural OT as a first
    production backend.  That would stack too many approximations before the
    underlying value and target contracts are audited.
  \item The most justified next HMC-relevant rung is EDH/PF with explicit
    importance weighting, preceded by EDH linear-Gaussian recovery.
  \item Soft and OT resampling should be studied explicitly as target-defining or
    target-altering relaxations, not as invisible implementation details.
  \item For nonlinear DSGE and MacroFinance models, DPFs should be benchmarked
    against existing exact or controlled-approximate BayesFilter paths before
    they are trusted as production HMC backends.
\end{enumerate}

This is a positive but guarded conclusion.  The literature and the student work
say the path is promising.  BayesFilter's doctrine adds the stronger claim that
promising is not enough: the project still needs per-rung value, gradient,
variance, and compiled-path evidence before a differentiable particle filter can
be called an HMC-safe backend for the structural model classes that matter here.
